\chapter{Literature Review}
\label{chap:literature}
\begingroup
\justifying
\setlength{\parindent}{2em}
\setlength{\parskip}{0.55\baselineskip}

\section{Why benchmark design matters}

A robot benchmark is sometimes treated as a collection of tasks plus a leaderboard. That description misses its scientific role. A useful benchmark specifies what changes, what remains fixed, how success is measured and which conclusions are licensed by the comparison. In robot learning, those choices are unusually consequential because the observation distribution, embodiment, controller and physical reset procedure can all change the apparent difficulty of a task.

Real-world evaluation has the strongest claim to operational relevance, but repeated trials are expensive and difficult to standardise. Objects wear, cameras move, human resetters differ and a failed manipulation can alter the next episode. Simulation trades some physical realism for repeatability, privileged state and rapid reset. It also makes counterfactual evaluation possible: the same policy can be tested with a new floor texture, a different store plan or a changed product split while other variables remain controlled. The trade is worthwhile only if the simulator preserves the geometric and interaction constraints that define the target domain.

Robot-learning suites occupy different points on this design spectrum. Meta-World and robosuite provide diverse tabletop tasks and modular control settings, while robomimic adds a systematic study of learning from offline human demonstrations~\cite{yu2020metaworld,zhu2020robosuite,mandlekar2021whatmatters}. ManiSkill2 broadens the range of embodiments, object models and manipulation families~\cite{gu2023maniskill2}. At room scale, iGibson 2.0 and Habitat 2.0 combine mobile manipulation with interactive household scenes and articulated objects~\cite{li2022igibson2,szot2021habitat2}. HomeRobot makes the integration problem explicit by evaluating navigation, novel-object grasping and placement in one open-vocabulary task~\cite{yenamandra2023homerobot}; ProcTHOR addresses a different bottleneck by generating large numbers of interactive training environments~\cite{deitke2022procthor}. Together, these systems show that task diversity, physical interaction and controlled generalisation can coexist. Their dominant settings are nevertheless tabletops and homes, not densely stocked stores.

Simulation also leaves a transfer question. RoboTHOR pairs simulated rooms with physical counterparts so that the drop between the two can be measured under aligned conditions, although its initial benchmark concerns embodied visual navigation rather than contact-rich retail manipulation~\cite{deitke2020robothor}. This distinction matters here: repeatable simulation is useful evidence, but it is not itself evidence of physical deployment.

The bias towards homes and tabletops is not superficial. Retail shelves contain repeated instances, near-duplicate packaging, tight clearances and fragile scene state. Fixtures create narrow aisles and long navigation distances. A policy may need to identify ``the Fanta bottle'' among visually adjacent branded products and then approach it from a pose compatible with the gripper. Household benchmark coverage therefore does not automatically imply retail competence. A dedicated benchmark is justified when the domain changes the distribution of geometry, appearance and failure, not simply when it supplies a new background texture.

\section{Retail datasets and the missing interaction layer}

Retail computer vision has produced substantial resources, but the datasets pose different problems. RPC records multi-product checkout scenes, whereas Products-10K concentrates on recognition among 10,000 fine-grained stock-keeping units~\cite{wei2019rpc,bai2020products10k}. The Grocery Store Dataset uses natural in-store images and a hierarchical label structure; the Freiburg Groceries Dataset supplies real-world grocery images at a coarser category level~\cite{klasson2019grocery,jund2016freiburg}. Across these resources, subtle packaging differences, clutter and domain shift make perception difficult. Their annotations still stop before physical action: they cannot determine whether a selected instance is reachable, whether the gripper collides with neighbouring stock or whether the object remains stable after placement.

Three-dimensional assets narrow part of this gap. Generic repositories such as YCB and ShapeNet have long supported manipulation and shape research~\cite{calli2015ycb,chang2015shapenet}; Google Scanned Objects adds more than a thousand simulation-ready scans of everyday items~\cite{francis2022gso}. Retail-specific collections better match product geometry and texture: IPA-3D1K contains reconstructed retail objects, while 3D-FRONT illustrates how textured assets and professionally designed layouts can be packaged as complete synthetic interiors~\cite{lindermayr2023ipa,fu2021threeDfront}. Access, licensing and scale remain practical constraints, and even a large mesh collection does not define a store. Assets still need consistent units, orientation, collision geometry and materials; they must then be placed on support surfaces, grouped like stock and simplified enough for interactive rendering.

Recent work begins to treat shelf picking as a generative robotics problem. FetchBot focuses on highly cluttered shelf arrangements, safety-aware retrieval and zero-shot transfer to a physical robot~\cite{liu2025fetchbot}. That focus is valuable, but an isolated shelf does not test movement through a store, variation in fixture layout or interaction with articulated refrigeration.

The closest neighbouring systems make the boundary sharper. FetchBench procedurally generates constrained fetching scenes---shelves, drawers, cabinets and baskets among them---and evaluates the full approach--grasp--retrieve chain with both modular and learned baselines~\cite{pmlr-v270-han25a}. It is deliberately object- and geometry-centred rather than retail-specific. Sari Sandbox moves in the other direction. It provides three virtual grocery-store layouts, more than 250 interactive products, an API-controlled embodied agent and SariBench human demonstrations for shopping tasks~\cite{gajo2025sari}. MarketGen generates complete supermarket environments from text or reference images, reports more than 1,100 goods and evaluates checkout unloading and in-aisle mobile manipulation with a sim-to-real study~\cite{hu2025marketgen}. RoboBenchMart, the benchmark used in this dissertation, pairs procedural dark-store layouts, stocked fixtures and motion-planned demonstrations in ManiSkill3, and reports fine-tuned VLA baselines under four evaluation conditions~\cite{soshin2025robobenchmart}. \ref{tab:retail-benchmark-comparison} summarises these systems without forcing them into a single ranking, and adds a row for the present study so that its narrower role is visible.

{\renewcommand{\arraystretch}{1.14}
\begin{table}[H]
\centering
\footnotesize
\begin{tabularx}{\textwidth}{
>{\RaggedRight\arraybackslash}p{0.16\textwidth}
>{\RaggedRight\arraybackslash}p{0.21\textwidth}
>{\RaggedRight\arraybackslash}X
>{\RaggedRight\arraybackslash}X}
\toprule
Resource & Scene and embodiment & Generated resources or data & Main evaluation emphasis \\
\midrule
FetchBench & Everyday fetching with a fixed-base Franka arm & Procedural shelves, drawers, cabinets and baskets; successful expert trajectories & Approaching, grasping and retrieving unknown objects in unseen constrained scenes \\
FetchBot & Densely cluttered shelves with a physical arm & Voxel-based simulated shelf scenes and an oracle-to-vision policy pipeline & Safety-aware retrieval and zero-shot sim-to-real transfer \\
Sari Sandbox / SariBench & Whole grocery stores with a VR avatar or API-driven agent & Three store layouts, 250+ interactive products and annotated human demonstrations & Navigation, inspection, manipulation, checkout and comparison with human behaviour \\
MarketGen & Complete supermarkets with cashier and mobile-manipulation agents & Agent-based procedural generation, 1,100+ goods and parameterised fixtures & Checkout unloading, in-aisle item collection and sim-to-real transfer \\
RoboBenchMart (upstream) & Dark-store retail with a Fetch mobile manipulator in ManiSkill3 & Procedural layouts, stocked fixtures, motion-planned demonstrations, fine-tuned checkpoints & Four VLA baselines under four conditions; atomic and composite tasks \\
This study & Same benchmark, current public release & No new data; 2,520 evaluation episodes with logs, videos and hashes & Independent re-evaluation of two released checkpoints and of the protocol itself \\
\bottomrule
\end{tabularx}
\caption{Scope comparison with adjacent retail and fetching systems. Counts and capabilities describe the cited releases; they are not normalised measures of benchmark scale or difficulty.}
\label{tab:retail-benchmark-comparison}
\end{table}}

The literature therefore suggests a layered gap. There is abundant 2D appearance data, a smaller number of 3D retail assets and even fewer systems that turn those assets into reproducible manipulation episodes. The missing layer is interaction: a task API, controller, success predicate and dataset path that let researchers move from ``which product is visible?'' to ``can the robot safely complete the requested operation?'' RoboBenchMart is an attempt to supply that layer. Whether its reported results can be recovered by an outside user is the question this dissertation takes up.

\section{Procedural environments and structured variation}

Procedural generation is attractive because a finite hand-built scene can be memorised. Randomness alone, however, is not enough. ProcTHOR uses procedural rules to scale interactive environments, Infinigen Indoors couples a constraint language with generated indoor assets, and Holodeck turns textual specifications into spatial constraints over retrieved assets~\cite{deitke2022procthor,raistrick2024infinigenindoors,yang2024holodeck}. Their mechanisms differ, but the shared lesson is simple: scene count matters only when generation also preserves structure. Uniformly scattering fixtures across a floor creates collisions, blocked aisles and arrangements that do not resemble stores.

RoboBenchMart's fixture placement draws on tensor-field methods for procedural street modelling~\cite{chen2008interactive}. In that setting, a field encodes locally preferred orientations; curves follow the field to form coherent road networks. A store presents a related geometric problem. Shelves usually align in rows, follow boundaries and leave passages between them. A second-order tensor is particularly useful because its principal directions are orientation-like: a shelf aligned at $0$ degrees is equivalent to one aligned at $180$ degrees. This symmetry avoids an arbitrary forward direction.

Several alternatives were possible. A fixed catalogue of plans would guarantee plausibility but provide limited continuous variation. Pure rejection sampling would diversify pose but struggle to produce coherent rows. Grammar-based layouts could encode strong design knowledge, although every new fixture relation would expand the grammar. A learned generator might capture examples from real plans but would require an appropriate corpus and a validity layer. The tensor-field approach occupies a practical middle ground: boundary and obstacle geometry induce local directions, while rejection tests still enforce collision and aisle constraints. Chapter~\ref{chap:system} shows how the upstream implementation actually uses the field, which is more discrete than the formulation suggests.

Procedural shelf population raises a different set of constraints. Products sit on detected support surfaces, typically in aligned groups. The \texttt{scene\_synthesizer} library provides scene and support-surface operations of this kind~\cite{eppner2025scene}. Visual diversity must also be accompanied by computational discipline. High-resolution product meshes multiply quickly, and a few hundred products per store can dominate memory and rendering time. Remeshing methods such as QuadriFlow~\cite{huang2018quadriflow} and simpler primitive approximations are the usual remedies; RoboBenchMart reports using simplified meshes for all but the nearest shelf~\cite{soshin2025robobenchmart}. This is an engineering optimisation rather than a learning contribution, yet it changes which experiments are feasible.

\section{Generating robot demonstrations}

Robot policy adaptation requires trajectories. Teleoperation gives human-like behaviour and can handle complex contact, but collection time grows with task and scene diversity. Physical teleoperation also carries reset cost. In simulation, privileged object poses and collision models make automatic collection possible. The central trade-off is between scalability and behavioural naturalness.

Sampling-based motion planning supplies collision-aware paths without learning a policy~\cite{LaValle_2006}. RRT-Connect grows two trees, one from the start and one from the goal configuration, for single-query planning~\cite{kuffner2000rrt}; OMPL packages this and other sampling-based algorithms behind reusable planning interfaces~\cite{sucan2012ompl}. Recent GPU systems take a different route: cuRobo combines parallel geometric search, inverse kinematics and trajectory optimisation for fast collision-free motion generation~\cite{sundaralingam2023curobo}. Asymptotically optimal sampling methods offer stronger theoretical properties at additional computational cost~\cite{KaramanSamplingBased}. For demonstration collection, however, the best planner is not necessarily the one with the strongest asymptotic guarantee. It is the one that produces enough successful, varied executions under a clear acceptance test.

RoboBenchMart's collector is of this pragmatic kind. Task logic defines waypoints for approach, grasp, transport and placement; between them, the planner first tries a screw motion~\cite{murray2017mathematical} and falls back to RRT-Connect when that fails. Only complete, successful episodes are kept. Such a scheme is scalable relative to manual demonstration, and seeded generation makes it reproducible in principle. It does not discover task structure automatically: waypoint design transfers human knowledge into the collector, and retaining only successes changes the trajectory distribution. The demonstrations show completion, but they do not teach recovery from error. Methods such as MimicGen take another route, transforming a small number of human source demonstrations across scenes and object poses~\cite{mandlekar2023mimicgen}. Reinforcement learning can generate behaviour through reward optimisation~\cite{sutton2018reinforcement,haarnoja2017soft}, yet reward design, exploration and computation remain substantial for long-horizon mobile manipulation. These families are complementary rather than mutually exclusive.

\section{Generalist vision-language-action policies}

Large heterogeneous robot datasets have encouraged policies that share representations across tasks and embodiments. Open X-Embodiment and DROID exemplify efforts to aggregate broad real-robot experience~\cite{open_x_embodiment_rt_x_2023,khazatsky2024droid}, addressing a data regime in which no single laboratory can cheaply collect every task and platform. Generalist policies then condition action generation on images, language and robot state.

Octo is a transformer-based generalist policy designed for adaptation across robots and tasks~\cite{octo_2023}. OpenVLA uses a vision-language backbone with autoregressive action prediction~\cite{kim24openvla}, Dita scales a diffusion transformer to the same problem~\cite{hou2025dita}, and the $\pi_0$ family pairs a vision-language foundation with a flow-based action component~\cite{black2024pi0visionlanguageactionflowmodel}. SmolVLA targets a smaller and more accessible model footprint~\cite{shukor2025smolvlavisionlanguageactionmodelaffordable}. $\pi_{0.5}$ extends the $\pi$ family toward open-world generalisation~\cite{intelligence2025pi05visionlanguageactionmodelopenworld}. These designs differ in scale, pretraining mixtures, action representation and adaptation recipe. A benchmark comparison therefore measures complete model pipelines, not model size in isolation.

Most of these models predict actions in chunks: a single inference returns several future actions, which are then executed before the policy is queried again. Chunking makes inference affordable and can smooth behaviour, but it also means that, for the length of a chunk, the policy cannot react to what it sees. How long a chunk is executed open loop is usually a deployment choice rather than a property of the model, and it is often reported only in passing.

Retail evaluation adds pressures that are less visible in a spacious tabletop setup. Grounding must discriminate branded instances in clutter. Manipulation needs centimetre-level alignment in constrained shelf cells. Mobile-base error changes the arm's reachable set. And when a task chains several stages, the probability of failure compounds.

The RoboBenchMart authors fully fine-tuned four of these models (Octo, SmolVLA, $\pi_0$ and $\pi_{0.5}$) on their demonstrations and released checkpoints for some of them~\cite{soshin2025robobenchmart}. Fine-tuning rather than zero-shot evaluation is appropriate for a domain-adaptation question, but it means each reported number belongs to a whole pipeline: pretrained weights, fine-tuning recipe, serving code and evaluation harness. When an outside user evaluates a released checkpoint, the first two parts are fixed and the last two are whatever the user's environment provides.

\section{Reproducibility of evaluation results}

Reproducibility problems in machine learning are well documented. In deep reinforcement learning, Henderson et al. showed that random seeds, implementation details and evaluation choices can change reported results by more than the differences between methods being compared~\cite{henderson2018deep}. The NeurIPS reproducibility programme responded with checklists and code-submission policies, while noting that sharing code is necessary but not sufficient~\cite{pineau2021improving}. Robot learning inherits these problems and adds its own: physical evaluation is costly to repeat, and simulated evaluation depends on simulator versions, assets and controller details that change over time.

Evaluating released policies in simulation has itself become a research topic. SimplerEnv, for example, builds simulated counterparts of real evaluation setups and checks whether policy rankings in simulation agree with those on hardware, paying explicit attention to control and visual mismatches between the two~\cite{li2024simplerenv}. The underlying lesson is that a policy checkpoint and an evaluation environment have to be matched carefully before their combination says anything about the policy.

Benchmark papers usually report results from a single evaluation infrastructure, maintained by the authors. Independent re-evaluations of the same released artefacts are rarer. I suspect one reason is that when the numbers do not match, the cause is hard to pin down after the fact, and an unexplained discrepancy is an awkward thing to publish. A re-evaluation is most useful when it records exactly what was run, so that a discrepancy can later be traced to a version, a configuration or a protocol detail rather than left as an unexplained gap.

\section{Generalisation, metrics and diagnostic evaluation}

Generalisation is meaningful only after its axis is named. A new robot pose tests local control robustness. A new texture tests visual invariance. A new floor plan changes navigation and camera geometry. An item that appeared in another task but not the current one tests whether the policy learned a transferable skill--object relation. A completely unseen object tests a stronger form of semantic and geometric transfer. Collapsing these cases into a single ``unseen'' score would hide useful structure.

The RoboBenchMart protocol defines a progression of this kind. Training-seed evaluation replays the configurations used for demonstrations; in-domain evaluation randomises the initial robot pose; unseen-scene evaluation changes the store; and unseen-scene-and-item evaluation pairs a familiar skill with an item not demonstrated for that task. The success rate is the primary metric,
\begin{equation}
    \widehat{S}_{m,t,c}=\frac{1}{N_{t,c}}\sum_{i=1}^{N_{t,c}} \mathbb{I}\!\left[\text{episode } i \text{ satisfies the task predicate}\right],
    \label{eq:success-rate}
\end{equation}
where $m$ denotes a model, $t$ a task, $c$ a condition and $N_{t,c}$ the number of episodes aggregated for that task and condition. The upstream study uses 100 episodes per task--item--fixture triplet. At that sample size, and more so at the 30 used here, the uncertainty in each rate is not negligible, and intervals rather than point estimates are the honest summary.

Success rates say how often a task is completed, not why it is not. The upstream authors address this with an eleven-category taxonomy that assigns each failed episode to the earliest stage at which it goes wrong. Choosing the earliest failure avoids double-counting cascades, at the price of hiding downstream damage, and the labels remain human judgements. A taxonomy of this kind needs per-episode evidence, ideally both video and state trajectories, to be applied consistently.

\section{Position of this dissertation}

This dissertation does not propose a model, a benchmark or a new evaluation condition. It takes one existing benchmark and two of its released checkpoints and asks whether the published results can be recovered from the public artefacts. Its contribution sits in the reproducibility layer described above: a complete, hashed record of what was run, a set of repairs needed to run it, and a careful reading of what the protocol does and does not measure. The benchmark design itself, including everything in Chapter~\ref{chap:system}, belongs to the RoboBenchMart authors.

\endgroup
